%%%%%%%%%%%%%%%%%%%%%%%%%%%%%%%%%%%%%%%%%%%%%%%%%%%%%%%%%%%%%%%%%%%%%%%%%%%%%
% AMBIENT COMPLEMENTARITY FROM THE CHAIN-LEVEL PERSPECTIVE
%%%%%%%%%%%%%%%%%%%%%%%%%%%%%%%%%%%%%%%%%%%%%%%%%%%%%%%%%%%%%%%%%%%%%%%%%%%%%

\subsection{Ambient complementarity from the chain-level perspective}
\label{sec:ambient-complementarity-chain-level}

Volume~I develops the ambient complementarity theorem
(Theorem~C and its nonlinear ambient-complementarity refinement)
at the level of cohomology and formal moduli.
The chain-level $\Ainf$ framework of the present volume gives
access to a more explicit model of the same geometry, in which
the $(-1)$-shifted symplectic structure and the complementary
Lagrangians are constructed directly from the bar complex and its
cyclic pairing.

\begin{definition}[Chain-level ambient complementarity complex;
\ClaimStatusConjectured]
\label{def:chain-level-ambient-complex}
Let $\cA$ be a chiral Koszul algebra with bar complex
$\barB_{\mathrm{ch}}(\cA)$
(Definition~\ref{def:chiral_bar}), Koszul dual $\cA^!$
with cobar complex $\Omegach(\barB_{\mathrm{ch}}(\cA))$
(Definition~\ref{def:chiral_cobar}), and universal
twisting morphism
$\tau_\cA \colon \barB_{\mathrm{ch}}(\cA) \to \cA^!$.
The \emph{chain-level ambient complementarity complex} is the
$\Ainf$ deformation complex
\[
T_{\mathrm{comp}}^{\Ainf}(\cA)
\;=\;
\mathrm{Def}_{\Ainf}(\cA)
\;\oplus\;
\mathrm{Def}_{\Ainf}(\cA^!)
\;\oplus\;
\mathrm{Def}_{\Ainf}(\tau_\cA),
\]
controlling the formal moduli problem of triples
$({\cA}', \alpha, {\cA^!}')$, where ${\cA}'$ is a curved
$\Ainf$ deformation of~$\cA$, ${\cA^!}'$ is a curved
$\Ainf$ deformation of~$\cA^!$, and $\alpha$ is an $\Ainf$
twisting morphism deforming $\tau_\cA$ compatibly with
${\cA}'$ and~${\cA^!}'$.  The differential on
$T_{\mathrm{comp}}^{\Ainf}(\cA)$ is the linearization of
the curved Maurer--Cartan equation
$d_{\barB}\alpha
+ \sum_{k \ge 2} m_k(\alpha^{\otimes k}) = 0$
for the full $\Ainf$ bar differential.
\end{definition}

\begin{remark}[The $(-1)$-shifted symplectic structure
from cyclic $\Ainf$ pairings]
\label{rem:shifted-symplectic-from-cyclic}
The cyclic structure on the bar complex
(Section~\ref{subsec:genus-tower-R}) endows
$T_{\mathrm{comp}}^{\Ainf}(\cA)$ with a degree~$-1$
pairing.  Concretely, the Verdier/BV trace on
$\barB_{\mathrm{ch}}(\cA)$ pairs bar elements of total
degree~$k$ with elements of degree $1-k$, so the induced
bilinear form on the deformation complex is antisymmetric
of degree~$-1$.  This is the chain-level avatar of the
$(-1)$-shifted symplectic structure on the formal moduli
problem $\mathcal{M}_{\mathrm{comp}}(\cA)$ constructed in
Volume~I.

The two summands $\mathrm{Def}_{\Ainf}(\cA)$ and
$\mathrm{Def}_{\Ainf}(\cA^!)$ are isotropic with respect
to this pairing: the Maurer--Cartan equation for a
one-sided deformation forces the self-pairing to vanish
(expanding the MC equation $d\gamma + \tfrac{1}{2}[\gamma,\gamma]
+ \cdots = 0$ and applying the cyclic pairing, the linear
term drops out by cyclicity and the quadratic term is
symmetric in the wrong degree).  Maximality of each
isotropic summand---i.e., the Lagrangian property---follows
from the bar-cobar quasi-isomorphism
(Theorem~\ref{thm:bar_cobar_adjunction}), which
identifies the normal complex to each Lagrangian with the
shifted cotangent fiber of the other.
\end{remark}

\begin{remark}[BV perspective: the complementarity
potential as a boundary master action]
\label{rem:bv-complementarity-potential}
In the holomorphic-topological framework
(Section~\ref{thm:boundary-chiral-algebra-bv}), the
complementarity potential $S_\cA$ of Volume~I admits a
chain-level interpretation as a \emph{boundary BV master
action}.  The ambient $(-1)$-shifted symplectic space
$\mathcal{M}_{\mathrm{comp}}(\cA)$ is the BV phase space
of the boundary theory; the two Lagrangians
$\mathcal{M}_\cA$ and $\mathcal{M}_{\cA^!}$ correspond to
bulk and boundary degrees of freedom respectively.  The
BV master equation
$\hh\,\Delta_{\mathrm{BV}} S_\cA
+ \tfrac{1}{2}\{S_\cA, S_\cA\} = 0$
is the formal-neighborhood form of the curved $\Ainf$
Maurer--Cartan equation on
$T_{\mathrm{comp}}^{\Ainf}(\cA)$, with the BV Laplacian
$\Delta_{\mathrm{BV}}$ corresponding to the genus-$g$
quantum corrections and the antibracket $\{-,-\}$
corresponding to the cyclic pairing.  The formal Legendre
duality between $S_\cA$ and $S_{\cA^!}$ (Volume~I,
Proposition~\ref*{prop:formal-legendre}) then becomes the
statement that bulk and boundary gauge-fixings are related
by fiber-base exchange in the shifted cotangent bundle
$T^*[-1]\mathcal{M}_\cA$.
\end{remark}

\begin{remark}[Nonlinear shadows and the quartic resonance class]
\label{rem:nonlinear-shadow-pointer}
The chain-level ambient complex
$T_{\mathrm{comp}}^{\Ainf}(\cA)$ gives in principle
access to all Taylor coefficients of the complementarity
potential $S_\cA$: the cubic coefficient is the universal
gravitational cubic coupling, and the quartic coefficient
governs the modular quartic resonance class
$\mathfrak{R}_{4,g,n}^{\mathrm{mod}}(\cA)$.  The full
nonlinear shadow calculus---including clutching laws,
family-by-family computations for Heisenberg, affine Kac--Moody,
and Virasoro, and the relationship between the quartic
resonance class and the universal modular master action
$\Theta_\cA$---is developed in Volume~I,
in the nonlinear modular shadows appendix.
\end{remark}


\subsection{The Kontsevich integral from the bar complex}
\label{subsec:kontsevich-from-bar}

The bar complex on $\FM_k(\C) \times \Conf_k(\R)$ computes the
Kontsevich integral of knots in $\C \times \R$.  This
identification connects the $\Ainf$ chiral algebra machinery to
concrete knot invariants and resolves Open Problem~1 of
Section~\ref{sec:conclusion} for the abelian case.

\subsubsection{Configuration space identification}

Recall the Kontsevich integral \cite{Kon94}.  Let
$K \colon S^1 \to \R^3$ be a Morse knot with respect to the
height function $t \colon \R^3 \to \R$.  The Kontsevich integral
is the formal sum
\begin{equation}
\label{eq:kontsevich-integral}
Z(K) \;=\; \sum_{n=0}^{\infty} \frac{1}{(2\pi i)^n}
\int_{t_{{\min}} < t_1 < \cdots < t_n < t_{{\max}}}
\sum_{P} \;\bigwedge_{(i,j) \in P}
\frac{d(z_i(t_i) - z_j(t_j))}{z_i(t_i) - z_j(t_j)}
\;\cdot\; D_P,
\end{equation}
where the sum is over pairings $P$ of the $n$ horizontal chords
connecting points on the knot, $z_i(t_i) \in \C$ is the
holomorphic coordinate of $K$ at height~$t_i$, and $D_P$ is
the chord diagram recording the pairing.  The integral is over
the ordered configuration $\Conf_n(\R)$ of heights.

Now consider the 3d holomorphic--topological theory on
$\C \times \R$ with Chern--Simons gauge group~$G$.  A knot
$K$ is a Wilson line operator: an element of the line category
$\mathcal{C}_{\mathrm{line}} \simeq \cA^!\text{-}\mathbf{mod}$
(Theorem~\ref{thm:lines_as_modules}).  The $\Ainf$ partition
function of this line operator is
\begin{equation}
\label{eq:ainfty-partition}
Z_{\Ainf}(K) \;=\; \sum_{n=0}^{\infty} \frac{1}{n!}
\int_{\Conf_n(\R)}
m_n\bigl(a_1(z_1,t_1), \ldots, a_n(z_n,t_n)\bigr)\,
dt_1 \wedge \cdots \wedge dt_n,
\end{equation}
where $m_n$ are the $\Ainf$ operations and the holomorphic
insertions $z_i = z_i(t_i)$ are determined by the knot~$K$.

The holomorphic weight forms entering~$m_n$
are integrated over $\FM_n(\C)$ (the Fulton--MacPherson
compactification of the holomorphic insertion points), and the
topological ordering is integrated over $\Conf_n(\R)$ (ordered
heights along the knot).  The integrand of $m_n$, restricted to
the cogenerator projection of the bar differential
(Theorem~\ref{thm:mc5-genus-zero-bridge}), is the logarithmic
form
\[
\omega_n \;=\; \bigwedge_{(i,j) \in E(T)} d\log(z_i - z_j)
\]
on $\FM_n(\C)$, summed over planar trees~$T$.  This is
\emph{exactly} the Kontsevich integrand
$\bigwedge d\log(z_i(t_i) - z_j(t_j))$ of
\eqref{eq:kontsevich-integral}.

\begin{proposition}[Configuration space identification;
\ClaimStatusProvedHere]
\label{prop:config-space-id}
The operation space $\FM_n(\C) \times \Conf_n(\R)$ of the
Swiss-cheese operad $\SCchtop$, restricted to configurations
along a Morse knot $K \colon S^1 \to \C \times \R$, is
canonically identified with the Kontsevich configuration space
$\{t_1 < \cdots < t_n\} \times \{z_i = z_i(t_i)\}_{i=1}^n$.
Under this identification, the bar differential's weight form
$\omega_n$ restricts to the Kontsevich integrand.
\end{proposition}

\begin{proof}
The identification is tautological: the Kontsevich
integral uses ordered heights $t_1 < \cdots < t_n$ along the
knot (an element of $\Conf_n(\R)$) and holomorphic positions
$z_1 = z_1(t_1), \ldots, z_n = z_n(t_n)$ in the transverse
plane (an element of $\C^n$, compactified to $\FM_n(\C)$ to
handle collisions).  The product
$\FM_n(\C) \times \Conf_n(\R)$ is the operation space
of~$\SCchtop$.  The weight form on $\FM_n(\C)$ is the
logarithmic form $\omega_n = \bigwedge_{(i,j)} d\log(z_i - z_j)$,
which restricts to $\bigwedge d\log(z_i(t_i) - z_j(t_j))$ along
the knot.
\end{proof}

\subsubsection{From the bar complex to knot invariants}

\begin{theorem}[Bar complex computes the Kontsevich integral;
\ClaimStatusConjectured]
\label{thm:bar-kontsevich}
Let\/ $\cA$ be a logarithmic\/ $\SCchtop$-algebra.  Let $\cA$ be the $\Ainf$ chiral
algebra of $G$-Chern--Simons theory on $\C \times \R$ at
level~$k$, and let $K \colon S^1 \to \C \times \R$ be a Morse
knot viewed as a line operator.  Then the $\Ainf$ partition
function \eqref{eq:ainfty-partition} equals the Kontsevich
integral \eqref{eq:kontsevich-integral}:
\[
Z_{\Ainf}(K) \;=\; Z(K)
\quad \in \quad \widehat{\cA}(\mathfrak{g}),
\]
where $\widehat{\cA}(\mathfrak{g})$ is the completed algebra of
chord diagrams valued in the Lie algebra~$\mathfrak{g}$.
In particular, evaluating $Z(K)$ via the $\mathfrak{sl}_2$
weight system at $q = e^{2\pi i/(k+2)}$ recovers the colored
Jones polynomial~$J_K(q)$.
\end{theorem}

\begin{remark}[Status and scope]
\label{rem:bar-kontsevich-status}
The configuration space identification
(Proposition~\ref{prop:config-space-id}) is proved.  The full
identification $Z_{\Ainf} = Z$ requires verifying that the
combinatorial coefficients---signs, symmetry factors, and
normalizations---match at all orders.  At genus~$0$ and for the
abelian case $G = \mathrm{U}(1)$, this is verified below.  The
nonabelian case requires the full MC5 identification
(Theorem~\ref{thm:mc5-genus-zero-bridge}) applied to
$\widehat{\mathfrak{g}}_k$, which is proved at all genera (Volume~I).
\end{remark}

\subsubsection{Verification: the abelian unknot}

For $G = \mathrm{U}(1)$ at level~$k$ and the unknot
$K = \{0\} \times \R$ (the trivial line operator along the
topological direction), the computation reduces to the Heisenberg
atom.

The $\Ainf$ operations are $m_2(J, J) = k \cdot \mathbf{1}$ and
$m_{n \geq 3} = 0$ (Section~\ref{sec:rosetta-stone}).  The
partition function \eqref{eq:ainfty-partition} has only the
$n = 0$ term (no insertions on the trivial knot contribute
nontrivially):
\[
Z_{\Ainf}(\text{unknot}) \;=\; \mathbf{1}
\;=\; \text{quantum dimension of the trivial representation}
\;=\; 1.
\]
The Kontsevich integral of the unknot is the empty chord diagram
$Z(\text{unknot}) = \mathbf{1}$ \cite{Kon94}, confirming the
identification.

For the unknot in the fundamental representation of $\mathrm{SU}(2)_k$,
the $R$-matrix $R(z)$ from Section~\ref{sec:spectral_braiding}
enters through the braiding of parallel line operators.  The
spectral $R$-matrix, composed $n$ times along the knot complement,
produces the quantum trace
\[
J_{\text{unknot}}(q) \;=\; \operatorname{tr}_V\bigl(
q^{2\rho}\bigr) \;=\; \frac{q - q^{-1}}{q^{1/(k+2)} -
q^{-1/(k+2)}} \;=\; [2]_q,
\]
the quantum dimension of the fundamental representation at
$q = e^{2\pi i/(k+2)}$, matching the standard result.

\begin{remark}[The trefoil and beyond]
\label{rem:trefoil}
For the trefoil knot $K_{3_1}$, the Kontsevich integral at order
$n = 1$ is
$Z_1(K_{3_1}) = \frac{1}{24}\,\theta$,
where $\theta$ is the unique chord diagram of degree~$1$.
In the bar complex, this corresponds to the single insertion
$m_2(J, J) = k$ integrated over the single crossing of the
trefoil's projection.  The coefficient $1/24$ matches the
Dedekind $\eta$-function contribution
$F_1 = -k/24 \cdot \log\,\eta(\tau)$ from
Section~\ref{sec:rosetta-stone}---the genus-$1$ modular
characteristic $\kappa(\cH_k) = k$ controls both the knot
invariant and the curved bar structure.  This coincidence is not
accidental: the trefoil's first-order Kontsevich coefficient and
the genus-$1$ free energy are both controlled by the same
Arnold defect on the elliptic curve.
\end{remark}

\begin{remark}[Relation to the Drinfeld--Kohno ladder]
\label{rem:dk-kontsevich}
The classical $r$-matrix
$r(z) = \int_0^\infty e^{-\lambda z}\,\{{\cdot}_\lambda{\cdot}\}
\,d\lambda$ produces the infinitesimal braiding, and the full quantum $R(z)$ produces the finite braiding.  For the affine lineage, this is proved unconditionally: one-loop exactness collapses the $A_\infty$ tower, the reduced HT monodromy is the KZ monodromy, and the Drinfeld--Kohno theorem identifies it with $U_q(\fg)$ (Theorem~\ref{thm:affine-monodromy-identification}).  On the $\mathfrak{sl}_2$ weight system, the bar complex integrals over $\FM_n(\C) \times \Conf_n(\R)$ compute the colored Jones polynomial (Corollary~\ref{cor:jones-polynomial}), recovering the Reshetikhin--Turaev invariant directly from the Swiss-cheese structure.  Extension beyond the evaluation locus to the full DK ladder (MC3) remains open in non-type-$A$ cases.
\end{remark}

\subsection{Braiding from the Steinberg involution}
\label{subsec:braiding-steinberg-involution}

\providecommand{\Steinberg}{\mathfrak{S}}
\providecommand{\Mvac}{\mathcal{M}_{\mathrm{vac}}}

The spectral $R$-matrix of Section~\ref{sec:spectral_braiding}
admits a geometric construction via the Steinberg variety of the
chiral algebra, in the spirit of Chriss--Ginzburg
\cite{ChrissGinzburg}.  The three axioms of the $R$-matrix---unitarity,
crossing symmetry, and the Yang--Baxter equation---become
consequences of the correspondence structure, rather than
independent verifications.

\subsubsection{The Steinberg involution}

\begin{construction}[The chiral Steinberg variety and its involution]
\label{constr:steinberg-involution}
Let $\cA$ be a chiral algebra with associated line category
$\mathcal{L}_b = \cA^!\text{-}\mathbf{mod}$
(Theorem~\ref{thm:lines_as_modules}) and vacuum moduli
$\Mvac$.  The \emph{chiral Steinberg variety} is the fiber
product
\[
\Steinberg_b \;=\; \mathcal{L}_b \times_{\Mvac} \mathcal{L}_b,
\]
parametrizing pairs of line operators that agree after
projection to the vacuum sector.  The \emph{Steinberg
involution}
\[
\sigma \colon \Steinberg_b \longrightarrow \Steinberg_b,
\qquad
(\ell_1, \ell_2) \longmapsto (\ell_2, \ell_1),
\]
exchanges the two factors of the fiber product.  This is the
geometric avatar of crossing symmetry: swapping the two line
operators that braid around each other corresponds to reversing
the crossing in the knot diagram.
\end{construction}

\subsubsection{The $R$-matrix axioms from the correspondence}

\begin{proposition}[Spectral $R$-matrix from the Steinberg involution;
\ClaimStatusProvedHere]
\label{prop:R-matrix-steinberg}
Let $R(z)$ be the spectral $R$-matrix of
Section~\ref{sec:spectral_braiding}, and let $\sigma$ be the
Steinberg involution of
Construction~\ref{constr:steinberg-involution}.  Define
$\rho = \tfrac{1}{2}\sum_{\alpha} \alpha$, half the sum of OPE
poles (the Weyl vector analog in the chiral setting).  Then:
\begin{enumerate}
\item \textbf{Unitarity.}
$R(z)\,R(-z) = \mathrm{id}$, as a consequence of
$\sigma^2 = \mathrm{id}$ on~$\Steinberg_b$.

\item \textbf{Crossing symmetry.}
$R(z)^{\sigma} = R(-z - \rho)$, from the interaction of the
involution~$\sigma$ with the $(-1)$-shifted symplectic form
on~$\Steinberg_b$.  The shift by~$\rho$ arises because the
symplectic form on the fiber product acquires a half-integer
twist from the OPE pole structure.

\item \textbf{Yang--Baxter equation.}
$R_{12}\,R_{13}\,R_{23} = R_{23}\,R_{13}\,R_{12}$, from
associativity of the triple convolution on the iterated
Steinberg
\[
\Steinberg_b^{(2)} \;=\;
\mathcal{L}_b \times_{\Mvac} \mathcal{L}_b
\times_{\Mvac} \mathcal{L}_b.
\]
\end{enumerate}
\end{proposition}

\begin{proof}
(1) The $R$-matrix is constructed as a correspondence class
$[R] \in H^*(\Steinberg_b)$ via convolution
(Section~\ref{sec:spectral_braiding}).  The composition
$R(z) \circ R(-z)$ is computed by the convolution product
on~$\Steinberg_b$, which factors through the diagonal
$\Delta \colon \mathcal{L}_b \hookrightarrow \Steinberg_b$.
The involution~$\sigma$ acts on~$[R]$ by
$\sigma^*[R(z)] = [R(-z)]$, so
$[R(z)] \star [R(-z)] = [R(z)] \star \sigma^*[R(z)]
= [\Delta_*\mathbf{1}] = \mathrm{id}$.

(2) The $(-1)$-shifted symplectic form~$\omega$ on the ambient
complementarity complex
(Remark~\ref{rem:shifted-symplectic-from-cyclic}) restricts
to~$\Steinberg_b$ as $\omega|_{\Steinberg_b}$.  Under the
involution, $\sigma^*\omega = -\omega + d\rho$, where the
exact correction~$d\rho$ arises from the half-sum of OPE
poles (the same mechanism that produces the $\rho$-shift in
the Harish-Chandra isomorphism).  Exponentiating:
$\sigma^*[R(z)] = [R(-z-\rho)]$.

(3) The triple convolution on~$\Steinberg_b^{(2)}$ is
associative because fiber products compose associatively.
The three compositions $R_{12} \star R_{13} \star R_{23}$ and
$R_{23} \star R_{13} \star R_{12}$ both factor through the
same triple fiber product
$\mathcal{L}_b \times_{\Mvac} \mathcal{L}_b
\times_{\Mvac} \mathcal{L}_b$
via the two orderings of the three projections.  The equality
of the two compositions is the content of the associativity
isomorphism for the convolution algebra
$H^*(\Steinberg_b)$, which is the Hecke algebra.
\end{proof}

\begin{remark}[Beyond the evaluation locus: the Springer
resolution analog]
\label{rem:springer-resolution-chiral}
The spectral $R$-matrix $R(z)$ of
Proposition~\ref{prop:R-matrix-steinberg} is constructed on the
\emph{regular locus} of~$\Steinberg_b$, where the fiber product
is smooth and the convolution algebra is well-defined.  Extending
$R(z)$ beyond the evaluation locus---from DK-$0$ to the full
spectral $R$-matrix---corresponds to extending the Steinberg
variety from the regular locus to the full (possibly singular)
correspondence.  This requires a resolution of singularities
of~$\Steinberg_b$: the \emph{chiral Springer resolution}
\[
\widetilde{\Steinberg}_b \;\longrightarrow\; \Steinberg_b,
\]
analogous to the classical Springer resolution
$\widetilde{\mathcal{N}} \to \mathcal{N}$ that resolves the
nilpotent cone.  The resolved variety
$\widetilde{\Steinberg}_b$ carries a well-defined convolution
product at all spectral parameters, yielding the full
non-evaluation $R$-matrix.  On the evaluation locus itself, the identification is settled: the reduced HT monodromy equals the quantum group $R$-matrix for all simple Lie algebras (Theorem~\ref{thm:affine-monodromy-identification}).  What remains open is the geometric construction of the chiral Springer resolution $\widetilde{\Steinberg}_b \to \Steinberg_b$ beyond type~$A$, which would extend the identification from the evaluation locus to the full spectral $R$-matrix.
(The spectral Drinfeld class vanishes for all simple Lie algebras
by root multiplicity one; Theorem~\ref{thm:complete-strictification}.)
\end{remark}
